% TOP_PROBLEM: {"id": 30006763, "title": "Uniform Boundedness Conjecture for Rational Points on Curves", "category_id": 1, "category": {"id": 1, "name": "number_theory", "display_name": "Number Theory", "description": "Properties of integers, prime numbers, Diophantine equations.", "slug": "number-theory", "order_index": 1, "created_at": "2026-07-31T15:26:25.670Z"}, "status": "open"}
% FIRST_POSTED: 2026-09-27
% !TeX program = lualatex
% Compile twice: lualatex open_problems_500.tex
% All references and prose are embedded; no BibTeX database or external inputs.
\documentclass[11pt,a4paper]{article}
\usepackage[margin=24mm,headheight=14pt]{geometry}
\usepackage{fontspec}
\setmainfont{Latin Modern Roman}
\setsansfont{Latin Modern Sans}
\setmonofont{Latin Modern Mono}
\usepackage{amsmath,amssymb,mathtools}
\usepackage{unicode-math}
\setmathfont{Latin Modern Math}
\usepackage{microtype}
\usepackage{enumitem,xurl,needspace}
\usepackage[dvipsnames]{xcolor}
\usepackage{hyperref}
\hypersetup{unicode=true,colorlinks=true,linkcolor=MidnightBlue,urlcolor=MidnightBlue,
 pdftitle={500 Mathematical Problem Statements},pdfauthor={Source-annotated compilation},bookmarksnumbered=true}
\usepackage{bookmark}
\usepackage{tocloft}
\setlength{\cftsecnumwidth}{3em}
\cftsetpnumwidth{2.5em}
\cftsetrmarg{4em}
\usepackage{titlesec}
\titleformat{\section}{\Large\bfseries\raggedright}{\thesection}{0.7em}{}
\usepackage{fancyhdr}
\pagestyle{fancy}\fancyhf{}
\fancyhead[L]{\small\sffamily Mathematical problem statements}
\fancyhead[R]{\small Source edition 22}
\fancyfoot[C]{\thepage}
\setlength{\parindent}{0pt}
\setlength{\parskip}{5pt plus 1pt}
\setlength{\emergencystretch}{3em}
\setcounter{tocdepth}{1}
\setcounter{secnumdepth}{1}
\setlist[itemize]{leftmargin=3em,itemsep=5pt,topsep=4pt,parsep=0pt}
\allowdisplaybreaks
\newcommand{\N}{\mathbb{N}}
\newcommand{\Z}{\mathbb{Z}}
\newcommand{\Q}{\mathbb{Q}}
\newcommand{\R}{\mathbb{R}}
\newcommand{\C}{\mathbb{C}}
\newcommand{\F}{\mathbb{F}}
\newcommand{\A}{\mathbb{A}}
\newcommand{\PP}{\mathbb P}
\newcommand{\E}{\mathbb{E}}
\newcommand{\eps}{\varepsilon}
\newcommand{\dd}{\,\mathrm d}
\newcommand{\sref}[2]{\hyperlink{source:#1:#2}{[S#2]}}
\newcommand{\eref}[2]{\hyperlink{extra:#1:#2}{[E#2]}}
% Turn the next switch off for a shorter reading copy. The quotations preserve
% every exactTarget field, including qualifications omitted from a short summary.
\newif\ifshowcatalogue
\showcataloguetrue
\newcommand{\cataloguescope}[1]{\ifshowcatalogue
 \paragraph{Exact scope in the supplied catalogue.}
 \begin{quote}\small #1\end{quote}\fi}
% Compile twice with: lualatex 30006763.tex
\numberwithin{equation}{section}
\hypersetup{pdftitle={Research attempt -- problem 30006763},pdfauthor={Research notebook; original compilation retained}}
\fancyhead[L]{\small\sffamily Research notebook}
\fancyhead[R]{\small\sffamily Problem 30006763 | 27 September 2026}
\begin{document}
\setcounter{section}{0}
% ================================================================
% problemId: 30006763
% Canonical problem ID: 30006763
% exactTarget SHA-256: c856fed9c514d21c26bba96b719ad4728bb53c2af04c0a76a286e4d15b87c5af
\clearpage
\section*{\texorpdfstring{Uniform Boundedness Conjecture for Rational Points on Curves}{Uniform Boundedness Conjecture for Rational Points on Curves}}\label{30006763}
\subsection{Definitions and mathematical statement}
Fix a number field $F$ and an integer $g\ge2$. In the usual curve convention, $C/F$ is smooth, projective, and geometrically connected, and its genus is $g=\dim_F H^0(C,\Omega_C^1)$. The requested uniform bound is
\[
 \exists c(g,F)<\infty\quad\forall C/F\text{ of genus }g,
 \qquad\#C(F)\le c(g,F).
\]
For an affine smooth curve, genus refers to its smooth projective completion and its rational points form a subset of those of that completion. The constant must not depend on coefficients, heights, or the isomorphism class of $C$. The target does not ask for a formula or an effective algorithm for the bound.
\par\smallskip\textit{Formulation and concept references:} \sref{30006763}{1}.
\cataloguescope{For each integer g >= 2 and number field F, prove that there is a constant c(g,F) such that every smooth curve C of genus g over F has at most c(g,F) F-rational points.}
\subsection{Short English statement}
With the number field and genus fixed, is there one finite bound for the number of rational points on every curve of that genus?
% BEGIN RESEARCH ADDITION ATTEMPT-164
\subsection{Research attempt}

\paragraph{Current frontier.}
Uniform Mordell--Lang bounds for a curve embedded in its Jacobian have
the form $c_1(g)c_2(g)^r$, where $r$ is the rank of the subgroup being
intersected. They do not remove rank from the bound
\sref{30006763}{1}. Yu--Yuan--Zhou's February 2026 manuscript makes the
constants explicit, still with Mordell--Weil rank in the exponent
\eref{30006763}{1}, Theorems 1.1--1.3. This recent theorem statement is not
being certified by an independent audit of its 178-page proof here; in
any event its rank dependence does not settle the printed target.
The classical conditional route via Lang's conjecture is different from
an unconditional genus-and-field bound \eref{30006763}{2}. The investigation
below asks whether the geometry of short sums of curve points can avoid
counting the entire Mordell--Weil lattice.

\paragraph{Attack route.}
One route is to bound the ranks of all relevant Jacobians, then use the
known rank-dependent theorem. That would require another major unresolved
uniformity assertion and may be much stronger than needed. A second is
to combine good reductions to a small finite group; the difficulty is
lifting modular collisions to genuine divisor equivalence. A third is
to use the large number of distinct low-degree effective divisors on a
curve together with a height-packing argument. I pursue the third and
quantify exactly which nonuniform parameters remain.

\paragraph{Attempt: short-sum rigidity from gonality.}
Assume $C(F)\ne\varnothing$, choose $P_0\in C(F)$, and let
$i(P)=[P-P_0]\in J(F)$. This map is injective for genus at least one:
a principal divisor $P-Q$ with $P\ne Q$ would give a degree-one map to
$\PP^1$. For a finite subset $\mathcal P\subset C(F)$ of size $N$, put
$A=i(\mathcal P)$.

Let $\operatorname{gon}_{\overline F}(C)$ be the least degree of a
nonconstant morphism $C_{\overline F}\to\PP^1$. If
$1\le d<\operatorname{gon}_{\overline F}(C)$, then
\begin{equation}\label{30006763:sumset}
 |dA|=\binom{N+d-1}{d},
 \qquad dA=\{a_1+\cdots+a_d:a_i\in A\}.
\end{equation}
To prove it, two distinct effective degree-$d$ divisors supported on
$\mathcal P$ cannot be linearly equivalent: cancelling their common
part would produce the zero and pole divisors of a nonconstant rational
function of degree at most $d$. Thus the map from size-$d$ multisets of
$\mathcal P$ into $\operatorname{Pic}^d(C)$ is injective. Counting those
multisets proves \eqref{30006763:sumset}. Equality of $F$-rational Jacobian
points implies the geometric divisor equivalence used in this argument,
so no descent hypothesis is hidden here.

In particular, for a nonhyperelliptic curve the additive energy is
\begin{equation}\label{30006763:sidon-energy}
 \mathcal E(A):=\#\{(a,b,c,d)\in A^4:a+b=c+d\}=2N^2-N.
\end{equation}
There are $N$ diagonal sums with one ordered representation and
$\binom N2$ other sums with two ordered representations. This proves
\eqref{30006763:sidon-energy} exactly, not just up to an implicit constant.

\paragraph{Attempt: the hyperelliptic correction.}
The preceding degree-two assertion is false for hyperelliptic curves,
including every genus-two curve. Let $\iota$ be the hyperelliptic
involution. Suppose $\mathcal P$ contains $h$ complete unordered pairs
$\{P,\iota P\}$ with $P\ne\iota P$, and $w$ fixed points of $\iota$.
The unique degree-two pencil contributes
$t=2h+w$ ordered pairs to one divisor class. Every other effective
degree-two divisor still has its unique unordered representation.
Consequently
\begin{equation}\label{30006763:hyper-energy}
 \mathcal E(A)=2N^2-N+t^2-4h-w
             =2N^2-N+t^2-2t+w\le3N^2.
\end{equation}
For the inequality use $0\le w\le t\le N$. The uniqueness of the
hyperelliptic pencil follows because two distinct degree-two maps would
make $C$ birational to a curve of bidegree $(2,2)$ in
$\PP^1\times\PP^1$, whose arithmetic genus is one; if the product map
has degree two, the maps have the same pencil. Either alternative rules
out distinct pencils in genus at least two.

The formula includes a subtle multiplicity check. A nonfixed conjugate
pair previously contributed $2^2=4$ to the energy, while a fixed point
previously contributed $1$. Replacing these separate contributions by
the single $t^2$ gives exactly \eqref{30006763:hyper-energy}, not a correction
of merely $t^2-t$. The script checks the combinatorics on formal sets
$\{\pm e_1,\ldots,\pm e_h\}$; those are tests of the identity, not
constructed counterexamples or actual curves over a number field.

\paragraph{Attempt: combining divisor rigidity with a lattice bound.}
Choose a symmetric ample height on $J$, with canonical quadratic height
$\widehat h$. Let
\[
 r=\operatorname{rank}J(F),\quad
 \tau=|J(F)_{\mathrm{tors}}|,\quad
 H=\max_{a\in A}\widehat h(a).
\]
For $r>0$ let $\mu>0$ be the minimum positive squared norm in the
Mordell--Weil lattice $J(F)/J(F)_{\mathrm{tors}}$ with this height.
Positivity follows from finite generation and positive definiteness of
the canonical height pairing. The triangle inequality for its Euclidean
norm puts $dA$ in the ball of radius $d\sqrt H$ modulo torsion.

Distinct lattice points have distance at least $\sqrt\mu$.
The disjoint open balls of radius $\sqrt\mu/2$ around lattice points in
a ball of radius $d\sqrt H$ lie in a ball of radius
$d\sqrt H+\sqrt\mu/2$. Comparing $r$-dimensional volumes gives
\begin{equation}\label{30006763:packing}
 |dA|\le\tau\left(1+2d\sqrt{H/\mu}\right)^r.
\end{equation}
There are at most $\tau$ lifts of any lattice point, which accounts for
the torsion factor. Combining \eqref{30006763:sumset} and
\eqref{30006763:packing} yields the proved finite-set bound
\begin{equation}\label{30006763:derived-bound}
 \boxed{\displaystyle
 \binom{N+d-1}{d}\le
 \tau\left(1+2d\sqrt{H/\mu}\right)^r
 \quad\text{for }d<\operatorname{gon}_{\overline F}(C).}
\end{equation}
If $r=0$, the simpler statement $N\le\tau$ applies; no nonexistent
minimum positive lattice length is introduced. Empty $C(F)$ needs no
argument. The quantities $H,\mu$ and $r$ belong to the chosen curve and
height; \eqref{30006763:derived-bound} is not uniform in those quantities.

For fixed $d$ this gives a bound of scale
$\tau^{1/d}(1+2d\sqrt{H/\mu})^{r/d}$ up to a factor depending on $d$.
It suggests why using divisors rather than points may help a quantitative
argument, but $d$ is bounded by the gonality of a fixed-genus curve and
cannot absorb arbitrary rank. No improvement over the strongest known
uniform Mordell--Lang bounds is claimed.

\paragraph{Stress test: why reductions do not close the gap.}
Suppose $C$ and $J$ have good reduction at finitely many places and put
$G_0=\prod_v J(k_v)$. Distinct sums may have the same image in $G_0$.
The finite-group energy inequality
$\mathcal E(B)\ge |B|^4/|G_0|$ therefore cannot be combined with
\eqref{30006763:sidon-energy} for the original set without a collision-lifting
lemma. This is not a technicality: on the free part of $J(F)$, the
kernel of the map to any finite group contains $|G_0|J(F)$ and hence has
positive rank whenever $r>0$. Reduction information alone cannot imply
that a difference in this kernel is zero. Special geometry of curve
points could still control these differences, but that is an additional
arithmetic theorem, not a property of the reduction map.

Likewise, low additive energy by itself does not bound cardinality in a
rank-free ambient group. The sets $\{e_1,\ldots,e_N\}\subset\Z^N$ satisfy
\eqref{30006763:sidon-energy} for every $N$, and arbitrarily large finite
Sidon sets also exist in $\Z$ if no height restriction is imposed. Thus
the geometric rigidity calculation cannot alone eliminate either rank
or the height scale. A putative compactness argument over the moduli
space must also contend with its noncompact degenerations and unbounded
arithmetic heights; none is suppressed here.

If one assumes a bound $r\le R(g,F)$ for all these Jacobians, the
established rank-dependent theorem immediately supplies the requested
$c(g,F)$. That is a valid conditional route, not a theorem established
in this entry. The desired bound on \emph{curve points} need not logically
require a bound on all Jacobian ranks; treating that stronger conjecture
as necessary would also overstate the reduction.

\paragraph{Outcome.}
\textbf{Outcome: Exploratory approach with unresolved gap.}

The exact short-sum count, the hyperelliptic energy correction, and the
height-sensitive packing inequality are proved as notebook lemmas, with
no claim of literature priority. They identify what geometry contributes
and which arithmetic parameters it fails to control. A new argument
uniformly controlling rational points in high-rank or high-height
situations is still required. Neither a counterexample family with
unbounded numbers of $F$-points nor the genus-and-field-only bound has
been obtained.

% Keep the local source heading and initial references together.
\Needspace{12\baselineskip}
% END RESEARCH ADDITION ATTEMPT-164
\subsection{Sources}
\begin{itemize}\raggedright
\item[\textnormal{[S1]}]\hypertarget{source:30006763:1}{} Kuhne, The number of rational points on a curve of genus at least two, EMS Press, 2023.
\url{https://ems.press/content/book-chapter-files/33181}.
{\small\textit{Catalogue formulation source.}}
% BEGIN RESEARCH ADDITION SOURCES-164
\item[\textnormal{[E1]}]\hypertarget{extra:30006763:1}{} Jiawei Yu, Xinyi Yuan and Shengxuan Zhou, ``Quantitativity on the number of rational points in the Mordell conjecture,'' arXiv:2602.01820v1 (2 February 2026), 178 pages. Theorems 1.1--1.3.
\url{https://arxiv.org/pdf/2602.01820}.
{\small\textit{Explicit rank-dependent bounds in a recent manuscript; no independent audit of the complete proof is claimed.}}
\item[\textnormal{[E2]}]\hypertarget{extra:30006763:2}{} Lucia Caporaso, Joe Harris and Barry Mazur, ``Uniformity of rational points,'' \emph{J. Amer. Math. Soc.} 10 (1997), 1--35, Theorem 1.1. See also the authors' ``Corrections to \emph{Uniformity of rational points} and further comments'' (arXiv title), \emph{Tunisian J. Math.} 4 (2022), 183--201; arXiv:2012.14461, Introduction.
\url{https://doi.org/10.1090/S0894-0347-97-00195-1}.
\url{https://arxiv.org/pdf/2012.14461}.
{\small\textit{Conditional uniformity via Lang; the later correction concerns a stronger generic-bound formulation.}}
% END RESEARCH ADDITION SOURCES-164
\end{itemize}

\end{document}
